% critique.tex — f166v, vv. 337–379 : commentaire critique
% Articles textuels tirés du commentaire de l'auteur (latex/Aristophane-Thesmo Harmattan 2023 vv 1-641.tex),
% texte et notes conservés ; répartition détaillée dans repartition.md.
\cc{346}{ἢ καὶ δέχεται προδιδοῦσ᾽ ἑταίρα τὸν φίλον.}{Je reprends la correction de Fritzsche qui substitue avec bonheur \grec{ἑταίρᾳ} à \grec{ἑταίρα}. Le scholiaste précise d'ailleurs qu'il y a différentes versions dans les manuscrits : certaines avec le iota souscrit, d'autres sans : 
\citComm{γρ<άφεται> ἔν τισιν ἑταίρα καὶ ἔστιν ἕτερος νοῦς [...] ἐὰν μὲν ἄνευ τοῦ ῑ γρ<άφηται> ἑταίρα [...] ἐὰν δὲ σὺν τῷ ῑ}
La leçon sans iota est incohérente : en quoi une hétaïre, dont c'est le métier de recevoir des dons en échange de faveurs \italique{trahirait-elle} (\grec{προδιδοῦσα}) son amant ce faisant ? 
(Willems, sentant bien la difficulté, s'éloigne en traduisant : « à l'insu de son amant ».) 
Seule donc la première version est satisfaisante : il est question au vers précédent d'une vieille faisant des cadeaux pour s'attirer les faveurs d'un galant (thème qu'Aristophane développe dans le \italique{Ploutos}, vv. 959 sqq.) ; après la lubricité, le lucre : cette fois la vieille reçoit des cadeaux pour servir d'entremetteuse en livrant (προδιδοῦσ᾽) le soupirant (τὸν φίλον) à une femme irrégulière (ἑταίρᾳ). (Rôle de vieille procureuse que tiendra d'ailleurs Euripide à la fin de la pièce.)}
